\chapter{Discussion}\label{ch:discussion}

\section{Interpretation of Key Findings}

\paragraph{Adaptive difficulty is the robust result.} Across every simulated world, including one with
no category weaknesses at all, the difficulty-aware policy almost eliminates puzzles the player has
little chance of solving (from 28--45\% down to 1--3\%). It keeps most practice in the productive zone.
On the real logs, the assumption that difficulty could be pitched from the player's Chess.com rating
proved badly wrong (a predicted 45\% solve rate against 84.6\% observed). A model that learns the
player's puzzle ability online fixes this. These two facts support each other. Pitching difficulty
correctly is both the most reliable benefit and the one most directly checked against real data.

\paragraph{Weakness targeting is real, but small and fragile.} Adaptive category selection beats random
selection in every simulation. The advantage shrinks by about two-thirds when weaknesses are as subtle
as the real cohort suggests. The new policy's \emph{extra} targeting edge over the original bandit does
not survive realistic effect sizes. The pre-registered replication tells the same story: the original
bandit clearly beats random, but falls short of the targeting threshold it was held to.

\paragraph{Game histories reveal level, not weaknesses.} This is the most consequential finding. From
how a player plays, their rating band can be predicted at almost three times chance, and their rating
within about 200 points. Yet which tactical categories they will miss in future games is barely more
predictable than population base rates. Per-category reliability is near zero. Two explanations fit.
Either category-specific differences between players are genuinely small, compared with overall
strength and with how hard each motif is for everyone. Or the measurement is too noisy to reveal them:
about 80 critical positions per player, spread over many categories, labelled by a tagger with $\kappa
\approx 0.5$. Both lead to the same design conclusion. A player's games can supply only a weak prior on
\emph{which} motif to train. Personal weaknesses have to be learned online, from targeted puzzles, and
that is what the IRT learner does.

\paragraph{What the negative results mean for the original plan.} The portfolio's gap was puzzle
\emph{generation} conditioned on the learner's weaknesses. The cohort result bears directly on it. If
per-category weakness cannot be reliably measured from typical game histories, a generator conditioned
on a diagnosed weakness would largely be conditioned on noise. The change of direction towards
difficulty-aware selection, with weakness learned online, is therefore supported by evidence, and was
not only a response to limited compute.

\section{Relation to Prior Work}

The finding that level is detectable from play is consistent with behavioural stylometry, which
identifies individuals from their games~\cite{mcilroyyoung2021detecting}. It is also consistent with
Maia's finding that play varies systematically with rating~\cite{mcilroyyoung2020aligning}. Those
models use hundreds of games and deep sequence encoders; with only 18 hand-made features and 60 games
per player, level is still recoverable. By contrast, the weak per-category signal agrees with Hristova
et al.~\cite{hristova2014assessing}, who found that difficulty comes mainly from the depth of
calculation rather than recognising the motif. A taxonomy by motif may simply not be where players
differ most. The case for learning online rather than profiling up front matches the tutoring-bandit
literature~\cite{clement2015multi} and Elo-based adaptive systems~\cite{pelanek2016applications}, which
estimate learner and item parameters jointly from interactions. Nie et al.~\cite{nie2026discovering}
point the same way at a much larger scale: they learn puzzle value from 1.5 billion interactions, not
from game analysis. Feng et al.~\cite{feng2025generating} remain the reference for puzzle
\emph{creation}, which this work did not attempt.

\section{Limitations of the Current Approach}\label{sec:limitations}

\subsection{How sophisticated is the AI in this system? An honest assessment}

\begin{table}[H]
\centering
\small
\caption{The AI components, how they work, and how much they contribute.}\label{tab:ai-honest}
\begin{tabularx}{\linewidth}{p{0.24\linewidth}p{0.22\linewidth}p{0.1\linewidth}Y}
\toprule
\textbf{Component} & \textbf{Technique} & \textbf{Learned?} & \textbf{Honest assessment} \\
\midrule
IRT learner with Thompson Sampling & Online Bayesian logistic model, full-covariance Laplace update & Online, from every attempt & The most sophisticated component and the core contribution. Mathematically careful (reduces to Glicko; tested), but a known family of models (Elo, IRT, extended Kalman) applied to a new problem. \\
PuzzleNet & Multi-task neural network (3{,}329 features, 512--256 trunk, three heads), written and differentiated by hand in NumPy & Offline, from 5.5M Lichess puzzles & The only deep-learning component, and the largest single jump in measured quality: tactic labelling improves from $\kappa=0.52$ to \KAPPANN{} on the same held-out puzzles, and mined-puzzle difficulty acquires a calibrated estimate with an uncertainty attached. It is still a feed-forward network over an engineered representation, not a convolutional or attention model over raw boards, and its labels imitate another automatic tagger rather than human judgement. \\
Beta--Bernoulli bandit & Thompson Sampling & Online & Textbook; a few lines of code. \\
Rating-band classifier & Logistic regression, random forest & From 300 real players & Standard supervised learning. Solid result, simple models. \\
Weakness models & Empirical Bayes; random forest & From 300 real players & Classical statistics. Shows only a small advantage over base rates (Chapter~\ref{ch:results}). \\
Glicko-2 fitter & Bayesian rating system & From logs & A faithful re-implementation of a published algorithm. \\
Tactic tagger, miner gates & Hand-written rules over engine output & No & Symbolic AI. Measured and much improved, but engineered, not learned. \\
Stockfish & Alpha-beta search with an NNUE evaluator & Pre-trained (third party) & By far the strongest AI in the system, and not a contribution of this project. \\
Generative model / deep RL & --- & --- & \textbf{Still none.} No puzzle generator was trained and no policy was learned by reinforcement learning; the bandit and the IRT learner are Bayesian, not deep. \\
\bottomrule
\end{tabularx}
\end{table}

In summary, the AI here is \textbf{mostly classical and statistically principled, with one trained
neural component}. PuzzleNet is a genuine supervised deep-learning model: 1.86M parameters, a
hand-derived backward pass, a noise-aware heteroscedastic likelihood, and calibration fitted on
held-out data. It is also, honestly, a multilayer perceptron over a carefully built representation:
a 1990s architecture with modern training practice, not the convolutional or attention models that
the published entries for this very task used~\cite{schutt2024estimating,milosz2024transformers}.
The remaining sophistication lies in three places: the formulation (a joint Bayesian model of ability,
category offset and difficulty, used for decisions under uncertainty), the engineering of reliable
measurement (deterministic analysis, win-probability grading, a validated tagger), and above all the
evaluation. That includes a pre-registered replication reported with its failures, paired and
Holm-corrected simulations, held-out validation, prequential scoring, and a 300-player real-data test
that challenges the system's own premise. Model capacity is still not where most of the value lies: the network was built
because two \emph{measured} failures pointed at it, and its worth is argued in the same currency as
everything else, on held-out data and downstream. An examiner comparing this with the portfolio will
see that its most ambitious element, generative reinforcement learning, is absent; this chapter argues
that the evidence justifies the change of direction.

\subsection{Other limitations}

\begin{itemize}
  \item \textbf{No human evaluation.} All evidence of learning benefit comes from simulation. Five real
        users give 14\% power to detect a medium effect; 34 paired participants are needed.
  \item \textbf{Labels are agreement, not truth.} PuzzleNet raises agreement with Lichess themes
        to \KAPPANN{}, but it is trained to imitate Lichess's own automatic tagger. Neither tagger has
        been validated against a human expert, so a blind spot shared by both would be invisible here.
  \item \textbf{Difficulty is population difficulty.} PuzzleNet predicts how hard a puzzle is for the
        Lichess population, not for this player, and it was calibrated on Lichess puzzles rather than
        on puzzles mined from a player's own games, where no ground truth exists.
  \item \textbf{A non-contextual bandit.} Selection ignores context such as session fatigue, time of day
        and the recent sequence of outcomes.
  \item \textbf{Simple persistence.} JSON files are adequate for a prototype but not for many users.
\end{itemize}

\section{Implications for AI-Assisted Chess Training}

\begin{enumerate}
  \item \textbf{Calibrate difficulty from puzzle outcomes, not from game rating.} Game and puzzle
        rating scales differ substantially, and a few attempts are enough to correct the offset.
  \item \textbf{Be sceptical of ``your weaknesses from your games'' claims.} At typical history
        lengths, per-motif diagnosis is mostly noise. Products that present a precise weakness
        breakdown after a few dozen games are likely presenting sampling noise with confidence.
        Honest interfaces should show uncertainty; this system shows a rating $\pm$ RD per category.
  \item \textbf{Use population norms.} ``Sacrifices are hard for everyone'' is a large and reliable
        effect, while ``you are bad at sacrifices'' is small and unreliable. Personal feedback should be
        relative to the norm.
  \item \textbf{Log what you need to evaluate yourself.} Recording each decision's propensity and the
        pre-outcome prediction turns every future user into evidence for off-policy
        evaluation~\cite{nie2026discovering}.
\end{enumerate}
